Concevoir une ligne de transport circulaire demande de choisir les lieux desservis
et l'ordre dans lequel les relier. Desservir davantage de zones peut raccourcir
l'accès des usagers, mais allonge parfois la ligne. À l'inverse, concentrer les
stations peut produire un anneau court tout en éloignant certains points de leur
station d'accès. Les deux décisions ne peuvent donc pas être traitées comme des
objectifs entièrement indépendants.

Le problème Ring-Star représente cette interaction dans un cadre géométrique
simple. Un anneau relie une partie des points ; les autres lui sont rattachés par
des arcs d'affectation. Le schéma de couverture illustre cette structure, avec les
arêtes du cycle en trait plein et les affectations en pointillé. Il est emprunté à
une publication consacrée à une variante bi-objective du problème\cite{figure}.
Le présent projet retient une somme pondérée des deux coûts et fixe le nombre
de stations, conformément à l'énoncé de la SAÉ\cite{enonce}.

La problématique étudiée est la suivante :
\begin{quote}\itshape
Comment sélectionner un nombre fixé de stations et construire leur cycle pour
réduire conjointement la longueur de la ligne et les distances d'accès, et quel
compromis entre qualité, temps de calcul et garantie offrent les méthodes proposées ?
\end{quote}

Une solution exacte sert de référence lorsque le problème et le budget de calcul
le permettent. Des méthodes approchées produisent rapidement des solutions
réalisables et cherchent à améliorer leur coût. La comparaison doit prendre en
compte la variabilité des recherches, le travail effectivement réalisé et le
statut des résultats du solveur. Une valeur renvoyée par une formulation exacte
ne devient pas un optimum prouvé du seul fait de l'utilisation d'une PLNE.

Ce rapport présente les méthodes effectivement implémentées dans le projet et
les résultats enregistrés de sa campagne de démonstration. Il explique les choix
de représentation, les règles de calcul et les contrôles qui permettent de relier
une figure à une solution précise. Les tableaux sont calculés depuis les exports
archivés ; aucune nouvelle campagne de grande ampleur n'est supposée réalisée.

La première partie développe le cadrage, la formulation mathématique, la
NP-difficulté, la PLNE et les trois méthodes approchées. La seconde décrit le
protocole, interprète les résultats, approfondit un cas sur \code{att48} et examine
la qualité logicielle ainsi que les limites de l'étude. Un glossaire et des annexes
permettent de retrouver les notations, les commandes et l'ensemble des coûts.
